\section{対称性で何を簡単にするか}\label{節:定積分の手法選択}

ここまでの方法は，どれも新しい原始関数を増やすものではない。定積分を，値の等しい別の積分に書き換え，二つを組にする方法である。最初に区間を見て，次にその区間を保つ変換を考えると，計算の候補を絞れる。

\begin{center}
\begin{tabularx}{0.96\linewidth}{@{}lXX@{}}
\toprule
区間や式の特徴 & 試す変換 & 比べる部分\\
\midrule
$[-a,a]$ & $x\mapsto -x$ & 偶関数と奇関数の部分\\
$[a,b]$ & $x\mapsto a+b-x$ & 和が定数になるか，一次因子が消えるか\\
$\left[0,\frac{\pi}{2}\right]$ & $x\mapsto\frac{\pi}{2}-x$ & 正弦と余弦が交換されるか\\
長い区間と周期関数 & 整数周期だけ平行移動 & 完全な周期と残りの区間\\
$\left[\frac1b,b\right]$ と $\frac{\dd x}{x}$ & $x\mapsto\frac1x$ & 逆数にした関数との和\\
\bottomrule
\end{tabularx}
\end{center}

候補が見つかっても，対称性だけで最後まで解けるとは限らない。対称性によって厄介な因子を一つ除き，残った部分に置換積分や部分積分を使えばよい。

\begin{example}{反転の後で置換と部分分数を使う}
次の定積分を求めよ。
\[
I=\int_0^\pi\frac{x\sin x}{1+\sin^2x}\dd x.
\]
\end{example}

$\frac{\sin x}{1+\sin^2x}$ は $x$ を $\pi-x$ に変えても同じである。したがって，最初に $x$ を対称性で取り除く。
\[
I=\frac{\pi}{2}J,\qquad
J=\int_0^\pi\frac{\sin x}{1+\sin^2x}\dd x.
\]
$J$ では分子の $\sin x\dd x$ が $-\dd(\cos x)$ に対応する。$u=\cos x$ と置くと，$1+\sin^2x=2-u^2$ だから
\[
J=\int_{-1}^1\frac{1}{2-u^2}\dd u.
\]
この分母は $(\sqrt2-u)(\sqrt2+u)$ に分かれる。よって
\[
\frac{1}{2-u^2}
 =\frac{1}{2\sqrt2}\left(\frac{1}{\sqrt2-u}+\frac{1}{\sqrt2+u}\right),
\]
\[
J=\left[\frac{1}{2\sqrt2}
\log\frac{\sqrt2+u}{\sqrt2-u}\right]_{-1}^1
 =\frac{1}{\sqrt2}\log\frac{\sqrt2+1}{\sqrt2-1}
 =\sqrt2\log(\sqrt2+1).
\]
ここでは $-1\leq u\leq1$ で対数の分子・分母がともに正である。最後の等式では $(\sqrt2+1)(\sqrt2-1)=1$ を使った。したがって
\[
I=\frac{\pi}{\sqrt2}\log(\sqrt2+1).
\]

もし最初に部分積分を使うなら，$\frac{\sin x}{1+\sin^2x}$ の原始関数が必要になり，さらにその対数を積分する項が現れる。反転によって $x$ を除いたため，必要な計算が一度の置換と部分分数で済んだ。同じ式でも，不定積分では「微分の逆を作れるか」を見て，定積分ではそれに加えて「区間を保つ置換で二つの式を補い合わせられるか」を見るのである。

